%%%%%%%%%%%%%%%%%%%%%%%%%%%%%%%%%%%%%%%%%%%%%%%%%%%%%%%%%%%%%%%%%%%%%%%%%%%%%%%%%%%%%%%%%%%%%
%%%%%%%%    Modified ICML 2021 Template     %%%%%%%%%%%%%%%%%%%%%%%%%%%%%%%%%%%%%%%%%%%%%%%%
%%%%%%%%%%%%%%%%%%%%%%%%%%%%%%%%%%%%%%%%%%%%%%%%%%%%%%%%%%%%%%%%%%%%%%%%%%%%%%%%%%%%%%%%%%%%%

\documentclass{article}

\usepackage{microtype}
\usepackage{graphicx}
\usepackage{subfigure}
\usepackage{booktabs}
\usepackage{amsmath}
\usepackage{amssymb}
\usepackage{hyperref}

\newcommand{\theHalgorithm}{\arabic{algorithm}}

\usepackage[accepted]{icml2021}

\begin{document}

\twocolumn[
\icmltitle{Exploratory Experiments on CNNs, ViT, CLIP, and VAEs}

\begin{icmlauthorlist}
\icmlauthor{Ali Hasnain}{}
\end{icmlauthorlist}

\icmlkeywords{Machine Learning, Computer Vision, Deep Learning}

\vskip 0.3in
]

% ============================================================
% ABSTRACT
% ============================================================

\begin{abstract}

% Summarize the motivation, experiments, major findings,
% and key conclusions of the entire paper.

\end{abstract}


% ============================================================
% 1. CNNs
% ============================================================

\section{CNNs}

\subsection{Introduction}

% Introduce convolutional neural networks and their role
% in hierarchical visual representation learning.

\subsection{Methodology}

% Describe:
% - Dataset
% - CNN architecture
% - Training setup
% - Baseline
% - Modified architecture
% - Feature extraction procedure
% - Evaluation metrics

\subsection{Results}

% Include:
% - Training/validation results
% - Feature visualizations
% - Tables
% - t-SNE/UMAP plots
% - Comparisons between models

\begin{figure}[ht]
    \centering
    \includegraphics[width=\columnwidth]{training_curves.png}
    \caption{Training and validation loss of the fine-tuned ResNet-152
    model over the training epochs.}
    \label{fig:resnet152_loss}
\end{figure}

\begin{figure}[ht]
    \centering
    \includegraphics[width=\columnwidth]{CNN_part1_predictions.png}
    \caption{Example predictions from the fine-tuned ResNet-152 model on
    images from the evaluation set.}
    \label{fig:resnet152_inference}
\end{figure}

\subsection{UMAP and t-SNE}

\begin{figure*}[ht]
    \centering

    \subfigure[UMAP with skip connections]{
        \includegraphics[width=0.48\textwidth]{umap-skip.png}
    }
    \hfill
    \subfigure[UMAP without skip connections]{
        \includegraphics[width=0.48\textwidth]{umap-noskip.png}
    }

    \vspace{0.15cm}

    \subfigure[t-SNE with skip connections]{
        \includegraphics[width=0.48\textwidth]{tsne-skip.png}
    }
    \hfill
    \subfigure[t-SNE without skip connections]{
        \includegraphics[width=0.48\textwidth]{tsne-noskip.png}
    }

    \caption{Comparison of learned feature representations using UMAP
    and t-SNE for ResNet-152 with and without skip connections. Each
    visualization shows features extracted from early, middle, and late
    layers of the network.}
    \label{fig:resnet_representation_comparison}
\end{figure*}

\subsection{Discussion}

% Interpret the results.
% Explain why the observed representations emerged.
% Discuss the effect of architectural modifications,
% skip connections, depth, etc.


% ============================================================
% 2. VISION TRANSFORMERS
% ============================================================

\section{Vision Transformers (ViT)}

\subsection{Introduction}

% Introduce Vision Transformers and their difference
% from convolutional architectures.

\subsection{Methodology}

% Describe:
% - ViT architecture
% - Dataset
% - Pretrained model, if applicable
% - Image preprocessing
% - Attention extraction
% - Layer/head selection
% - Evaluation procedure

\begin{figure}[ht]
    \centering
    \includegraphics[width=\columnwidth]{imgs.png}
    \caption{Sample images used in the ViT experiments.}
    \label{fig:vit_images}
\end{figure}

\subsection{Results}

% Include:
% - Classification results
% - Attention maps
% - Attention visualizations
% - Intermediate representation analysis
% - Comparisons across layers/heads

\begin{figure}[ht]
    \centering
    \includegraphics[width=\columnwidth]{finalattn.png}
    \caption{Attention overlays from the final ViT layer, visualizing the image regions receiving the highest attention.}
    \label{fig:vit_final_attention}
\end{figure}

\begin{figure}[ht]
    \centering
    \includegraphics[width=\columnwidth]{attentionacrosslayers.png}
    \caption{Attention patterns across different ViT layers for a dog image.
    The top row shows CLS-to-patch attention averaged across all heads for
    layers 0, 2, 4, 6, 8, and 11, illustrating the evolution of spatial
    attention throughout the network. The bottom row shows attention from
    six individual heads in the final layer, highlighting the diversity of
    attention patterns across heads.}
    \label{fig:vit_attention_layers}
\end{figure}

\begin{figure}[ht]
    \centering
    \includegraphics[width=\columnwidth]{layer4attn.png}
    \caption{Attention overlays from Layer 4 of the ViT. The visualization
    shows the spatial regions receiving attention from the CLS token after
    averaging attention across all heads.}
    \label{fig:vit_layer4_attention}
\end{figure}

\begin{figure}[ht]
    \centering
    \includegraphics[width=\columnwidth]{maskingex.png}
    \caption{Effect of patch masking on a ViT input image. The original
    image is compared with inputs where 30\% of the image patches are
    masked either randomly or according to their proximity to the image
    center. Each masked patch is replaced with a zero-valued (black)
    region.}
    \label{fig:vit_masking}
\end{figure}

\begin{figure}[t]
    \centering
    \includegraphics[width=\columnwidth]{maskingconfidence.png}
    \caption{ViT confidence under increasing levels of patch masking.
    Each subplot corresponds to a different input image. Confidence is
    measured for the class predicted from the original unmasked image,
    while patches are progressively masked using either random or
    center-based masking.}
    \label{fig:vit_masking_confidence}
\end{figure}

\begin{figure}[ht]
    \centering
    \includegraphics[width=\columnwidth]{maskingprogression.png}
    \caption{Qualitative effect of progressively masking ViT input patches.
    For each input image, the top row shows random masking and the bottom
    row shows center masking at increasing masking fractions. The
    confidence values indicate the model's confidence in the class
    predicted for the original unmasked image.}
    \label{fig:vit_masking_progression}
\end{figure}

\begin{table}[t]
    \centering
    \caption{Effect of patch masking on ViT predictions. Values show the
    confidence in the original prediction class under random and center
    masking. Bold values indicate higher confidence between the two
    masking strategies.}
    \label{tab:masking_results}
    \small
    \begin{tabular}{c|cc|cc|cc}
        \toprule
        & \multicolumn{2}{c|}{\textbf{Golden Retriever}}
        & \multicolumn{2}{c|}{\textbf{Sports Car}}
        & \multicolumn{2}{c}{\textbf{Egyptian Cat}} \\
        \textbf{Mask} & Random & Center & Random & Center & Random & Center \\
        \midrule
        0\%  & 98.34 & 98.34 & 75.84 & 75.84 & 58.10 & 58.10 \\
        10\% & 98.57 & 97.55 & 85.34 & 75.96 & 59.40 & 70.74 \\
        20\% & 97.54 & 87.18 & 62.61 & 23.36 & 52.80 & 59.68 \\
        30\% & 97.98 & 61.63 & 78.69 & 23.92 & 55.35 & 58.68 \\
        40\% & 94.43 & 42.82 & 83.14 & 14.95 & 63.66 & 51.75 \\
        50\% & 97.29 & 47.89 & 77.39 & 20.19 & 39.95 & 60.85 \\
        60\% & 58.36 & 11.16 & 52.79 & 17.56 & 52.43 & 59.80 \\
        70\% & 94.77 & 3.53  & 61.00 & 17.08 & 48.81 & 48.13 \\
        80\% & 7.02  & 3.71  & 35.92 & 9.30  & 39.37 & 44.99 \\
        \bottomrule
    \end{tabular}
\end{table}

\begin{table}[ht]
    \centering
    \caption{Selected prediction changes under patch masking.}
    \label{tab:masking_predictions}
    \small
    \begin{tabular}{ccll}
        \toprule
        \textbf{Image} & \textbf{Mask} & \textbf{Random} & \textbf{Center} \\
        \midrule
        Golden Retriever & 50\% &
        Golden Retriever &
        Matchstick \\
        
        Golden Retriever & 60\% &
        Golden Retriever &
        Nematode \\
        
        Golden Retriever & 80\% &
        Panpipe &
        Panpipe \\
        
        Sports Car & 60\% &
        Convertible &
        Car mirror \\
        
        Sports Car & 70\% &
        Car mirror &
        Car mirror \\
        
        Egyptian Cat & 50\% &
        Tabby cat &
        Egyptian cat \\
        \bottomrule
    \end{tabular}
\end{table}

\begin{figure}[ht]
    \centering
    \includegraphics[width=\columnwidth]{clsmean.png}
    \caption{Comparison of linear probe performance using the final [CLS] token and mean-pooled patch token representations. Both representations achieve 100\% training accuracy, while the [CLS] representation achieves slightly higher test accuracy (95.67\%) than mean patch pooling (95.33\%).}
    \label{fig:cls-vs-mean}
\end{figure}

\subsection{Discussion}

% Discuss:
% - What attention learns
% - Spatial structure
% - Attention sinks
% - Differences between early/middle/late layers
% - Comparison with CNN representations


% ============================================================
% 3. CLIP
% ============================================================

\section{Contrastive Language-Image Pre-training (CLIP)}

\subsection{Introduction}

Contrastive Language--Image Pre-training (CLIP) is a multimodal
representation learning framework introduced by OpenAI that learns a
shared embedding space for images and natural language. Instead of
training directly on a fixed set of class labels, CLIP is pretrained
on large-scale image--text pairs using a contrastive objective. Given
an image and a collection of text descriptions, the model learns to
increase the similarity between the corresponding image--text pair
while decreasing the similarity to mismatched pairs.

This shared representation enables CLIP to perform zero-shot
classification without requiring task-specific classifier heads.
For example, an image can be compared against text prompts such as
``a photo of a dog'' or ``a photo of a cat'', with the class
corresponding to the most similar text embedding selected as the
prediction. This provides a flexible alternative to conventional
supervised image classification, where a model is trained explicitly
for a fixed set of classes.

The original CLIP study demonstrated that large-scale
vision-language pretraining can produce strong zero-shot transfer
performance across a wide range of datasets and tasks. Its ability to
associate visual concepts with natural language also makes CLIP useful
for applications where multiple modalities need to be represented
within a common semantic space. The model consists of separate image
and text encoders whose outputs are projected into the same embedding
space, allowing their representations to be directly compared.

In this experiment, we investigate the behavior of CLIP under
different prompting strategies for zero-shot classification. In
particular, we compare plain class labels, simple natural-language
templates, and more descriptive prompts. We further examine the
structure of the resulting image and text embedding spaces and
investigate whether Procrustes alignment can improve their
correspondence.

\subsection{Methodology}

% Describe:
% - CLIP model
% - Image and text encoders
% - Dataset
% - Embedding extraction
% - Similarity computation
% - Procrustes alignment, if applicable
% - Visualization procedure


\begin{table}[ht]
    \centering
    \caption{Prompt strategies with representative examples.}
    \label{tab:prompt_strategies}
    \small
    \begin{tabular}{lp{0.52\columnwidth}}
        \toprule
        \textbf{Strategy} & \textbf{Example Prompt} \\
        \midrule
        Plain labels
        & \texttt{airplane} \\[2pt]
        A photo of
        & \texttt{a photo of an airplane} \\[2pt]
        Descriptive
        & \texttt{a photo of an airplane in the sky} \\
        \bottomrule
    \end{tabular}
\end{table}

\subsection{Results}

% Include:
% - Image/text embedding visualizations
% - Before/after alignment
% - Class-level clustering
% - Similarity results
% - Tables and figures

\begin{table}[ht]
    \centering
    \caption{Zero-shot classification accuracy using different text prompt formulations.}
    \label{tab:prompt_results}
    \begin{tabular}{lc}
        \toprule
        \textbf{Prompt Type} & \textbf{Accuracy (\%)} \\
        \midrule
        Plain labels & 96.26 \\
        ``A photo of'' & \textbf{97.36} \\
        Descriptive & 90.80 \\
        \bottomrule
    \end{tabular}
\end{table}

\begin{table}[ht]
    \centering
    \caption{In-sample zero-shot classification accuracy before and after Procrustes alignment.}
    \label{tab:clip_alignment}
    \begin{tabular}{lc}
        \toprule
        \textbf{Method} & \textbf{Accuracy (\%)} \\
        \midrule
        Original CLIP & 97.36 \\
        After Procrustes & \textbf{100.00} \\
        \midrule
        Change & \textbf{+2.64 pp} \\
        \bottomrule
    \end{tabular}
\end{table}

\begin{figure}[t]
    \centering

    \subfigure[Before normalization]{
        \includegraphics[width=0.95\columnwidth]{anorm.png}
    }

    \vspace{0.15cm}

    \subfigure[After normalization]{
        \includegraphics[width=0.95\columnwidth]{bnorm.png}
    }

    \caption{Visualization of the embedding space before and after normalization.}
    \label{fig:embedding_normalization}
\end{figure}

\begin{figure}[t]
    \centering
    \includegraphics[width=\columnwidth]{clip_procrustes.png}
    \caption{Visualization of CLIP embeddings after Procrustes alignment.}
    \label{fig:clip_procrustes}
\end{figure}



\subsection{Discussion}

% Discuss:
% - Alignment between image and text spaces
% - Semantic structure
% - Effectiveness of Procrustes alignment
% - Limitations of the experiment


% ============================================================
% 4. VARIATIONAL AUTOENCODERS
% ============================================================

\section{Variational Autoencoders (VAE)}

\subsection{Introduction}

% Introduce VAEs and latent-variable representation learning.

\subsection{Methodology}

% Describe:
% - Encoder
% - Decoder
% - Latent dimensionality
% - Reparameterization trick
% - Reconstruction loss
% - KL-divergence
% - Training setup

\subsection{Results}

% Include:
% - Reconstruction examples
% - Training curves
% - Latent-space visualization
% - MNIST embeddings
% - Generated samples
% - Effect of latent dimensionality

\begin{figure}[ht]
    \centering
    \includegraphics[width=\columnwidth]{vae_training_loss.png}
    \caption{VAE training loss over training epochs. The loss captures the
    combined reconstruction and KL-divergence objectives used to optimize
    the variational autoencoder.}
    \label{fig:vae_training_loss}
\end{figure}

\begin{figure}[t]
    \centering
    \includegraphics[width=\columnwidth]{samples_generated_from_vae.png}
    \caption{Samples generated by the trained VAE by decoding latent
    representations sampled from the learned latent distribution.}
    \label{fig:vae_generated_samples}
\end{figure}

\begin{figure}[t]
    \centering
    \includegraphics[width=\columnwidth]{reconstruction.png}
    \caption{Original and reconstructed images produced by the trained
    VAE. The reconstructions illustrate how well the decoder preserves
    the salient visual features of the input images.}
    \label{fig:vae_reconstruction}
\end{figure}

\begin{figure}[t]
    \centering
    \includegraphics[width=\columnwidth]{latent.png}
    \caption{Two-dimensional latent space learned by the VAE on the MNIST
    test set. Each point represents the latent mean of a test image, with
    points colored according to their digit class.}
    \label{fig:vae_latent_space}
\end{figure}

\subsection{Discussion}

% Discuss:
% - Structure of the learned latent space
% - Class separation
% - Reconstruction quality
% - Trade-off between reconstruction and regularization
% - Limitations


% ============================================================
% 5. CONCLUSION
% ============================================================

\section{Conclusion}

% Summarize the key findings from all four experiments.
% Discuss what the experiments collectively reveal about
% representation learning across CNNs, ViTs, CLIP, and VAEs.
%
% Mention limitations and possible future directions.


% ============================================================
% REFERENCES
% ============================================================

\bibliography{example_paper}
\bibliographystyle{icml2021}

\end{document}